% Section content template for individual Educative lessons
% This file contains the actual content for a specific section
% Generated from Educative API section content

% Section: Evaluation of the Typeahead Suggestion System’s Design
% Section ID: 6626894497185792
% Section Slug: evaluation-of-the-typeahead-suggestion-systems-design
% Generated: 2025-12-08T15:11:39.259023

\subsection{Requirements compliance}\label{MQypNT8H9uVEDRNoSTQmr}
\\
The non-functional requirements of the proposed typeahead suggestion system are low latency, fault tolerance, and scalability.

\begin{itemize}
\item
\protect\phantomsection\label{kxAqZJuVqO8N2nyxxrz1c}
\textbf{Low latency:} There are various levels at which we can minimize the system's latency. We can minimize the latency with the following options:

\begin{itemize}
\item
Reduce the depth of the tree, which reduces the overall traversal time.
\item
Update the trie offline, which means that the time taken by the update operation isn't on the clients' critical path.
\item
Use geographically distributed applications and database servers. This way, the service is provided near the user, which also reduces any communication delays and aids in reducing latency.
\item
Use Redis as a caching layer on top of NoSQL database clusters and Cassandra database to ensure low-latency data retrieval and improved performance.
\item
Appropriately partition tries, which leads to a proper distribution of the load and results in better performance.
\end{itemize}
\item
\protect\phantomsection\label{jHGBEhKx6C_OSvxJgtr-a}
\textbf{Fault tolerance:} Since the replication and partitioning of the trees are provided, the system operates with high resilience. If one server fails, others are on standby to deliver the services.
\item
\protect\phantomsection\label{BCR9Hri_I-h1uMOFjiJlh}
\textbf{Scalability:} Since our proposed system is flexible, more servers can be added or removed as the load increases. For example, if the number of queries increases, the number of partitions or shards of the trees is increased accordingly.
\end{itemize}

\textbf{Non-functional Requirements Compliance}

\begin{table}[htbp]
\centering
\small
\adjustbox{max width=\textwidth}{
\begin{tabular}{|p{0.351\textwidth}|p{0.599\textwidth}|}
\hline
\textbf{Non-functional Requirements} & \textbf{Approaches} \\
\hline
\hline
\hfill\break Low latency & \begin{itemize} \tightlist \item Reducing the depth of the tries makes the traversal faster \item Updating the tries offline and not in real time) \item Partitioning of the tries \item Caching servers \end{itemize} \\
\hline
Fault tolerance & \begin{itemize} \tightlist \item Replicating the tries and the NoSQL databases \end{itemize} \\
\hline
Scalability & \begin{itemize} \tightlist \item Adding or removing application servers based on the incoming traffic \item Increasing the trie partitions \end{itemize} \\
\hline
\end{tabular}
}
\caption{Non-functional Requirements Compliance}
\end{table}

\subsection{Client-side optimisation}\label{1xq6fh2F1ES2gtMuyrY2I}
\\
To improve the user's experience, we can implement the following client-side optimizations:

\begin{itemize}
\item
\protect\phantomsection\label{TmsizmC0lYAyh4QEWoAJU}
The client should only attempt to contact the server if the user hasn't pressed any keys for some time---for example, any delay greater than 160 ms, which is the average delay between two keystrokes. This way, we can also avoid unnecessary bandwidth consumption. This suggestion might not be useful when a user is typing rapidly.
\item
\protect\phantomsection\label{DvKrwefRi6DF6NbyE-PXT}
The client can initially wait till the user types a few characters.
\item
\protect\phantomsection\label{tzC7cMy2pkb41VtzG2_Qv}
Clients can save a local copy of the recent history of suggestions. The rate of reuse of recent history in the suggestions list is relatively high.
\item
\protect\phantomsection\label{pW3ApC3PCJQHMdTZFsIJz}
One of the most crucial elements is establishing a connection with the server as soon as possible. The client can establish a connection with the server as soon as the user visits the search page. As a result, the client doesn't waste time establishing the connection when the user inputs the first character. Usually, the connection is established with the server via a \textbf{WebSocket protocol}.
\item
\protect\phantomsection\label{P5P8lx5N2YZxRyvC1dkW_}
For efficiency, the server can push a portion of its cache to CDNs and other edge caches at Internet exchange points (IXPs) or even inside a client's Internet service provider (ISP).
\end{itemize}

\subsection{Personalization}\label{djEzDBFCiubx23psYcPiU}
\\
Users receive typeahead suggestions based on their previous searches, location, language, and other factors. We can save each user's personal history on the server separately and cache it on the client. Before transmitting the final set to the user, the server might include these customized phrases. Personalized searches should always take precedence over other types of searches.

\subsection{Summary}\label{Lk0w74uLtk_sXVJdrHDC5}
\\
In this design problem, we learned how pushing resource-intensive processing to the offline infrastructure and using appropriate data structures enables us to serve our customers with low latency. Many optimizations lend themselves to specific use cases. We saw multiple optimizations on our trie data structures for condensed data storage and quick serving.

% End of section content